\documentclass[aspectratio=169]{beamer}

% ================================================================
% Building AI Agents for Production
% 90-minute technical session
% ================================================================

\usepackage{../../../shared/theme/esmad_beamer_theme}
\usepackage{listings}
\usepackage{booktabs}

% Presentation information
\title{Building AI Agents}
\subtitle{Architecture, Reliability, and Production Operations}
\date{\today}

\begin{document}

\begin{frame}
  \titlepage
\end{frame}

\tocslide

% ================================================================
\section{Session Plan}
% ================================================================

\begin{frame}{90-Minute Session Plan}
  \begin{table}
    \small
    \begin{tabular}{@{}ll@{}}
      \toprule
      Time & Block \\
      \midrule
      0-10 min & Agent definitions and decision criteria \\
      10-30 min & Architecture and orchestration patterns \\
      30-50 min & Tool contracts, memory, retrieval \\
      50-70 min & Safety, evaluation, and observability \\
      70-85 min & End-to-end case study walkthrough \\
      85-90 min & Final checklist and Q\&A \\
      \bottomrule
    \end{tabular}
  \end{table}
\end{frame}

\begin{frame}{Delivery Modes}
  \textbf{Use this same deck in three modes:}
  \begin{itemize}
    \item \textbf{Core 60 min:} skip slides marked as ``Deep Dive''
    \item \textbf{Standard 75 min:} include selected Deep Dive slides
    \item \textbf{Full 90 min:} run all slides with checkpoints and discussion
  \end{itemize}

  \vspace{0.5em}
  \begin{block}{Deep Dive marker}
    Slides labeled \textbf{Deep Dive} are optional expansion content.
  \end{block}
\end{frame}

\begin{frame}{Learning Outcomes}
  By the end of this session, you should be able to:
  \begin{itemize}
    \item Decide when to use an agent vs a deterministic workflow
    \item Design an agent architecture with explicit control boundaries
    \item Define robust tool contracts and memory policies
    \item Build an evaluation loop that measures quality and cost
    \item Identify failure modes before they reach production
  \end{itemize}
\end{frame}

% ================================================================
\section{Agent Fundamentals}
% ================================================================

\begin{frame}{What Is an AI Agent?}
  \begin{definitionbox}{Production Agent}
    A software system that iteratively selects actions, uses tools, updates state, and pursues a task objective under explicit safety, cost, and termination constraints.
  \end{definitionbox}

  \vspace{0.4em}
  \textbf{Properties that matter:}
  \begin{itemize}
    \item Explicit objective and termination criteria
    \item Observable state transitions
    \item Typed integration surfaces
    \item Safety and governance controls
  \end{itemize}
\end{frame}

\begin{frame}{Agent vs Single LLM Call}
  \begin{columns}
    \column{0.5\textwidth}
    \textbf{Single LLM call}
    \begin{itemize}
      \item One-shot response
      \item No execution state
      \item Limited external action
      \item Best for simple transformations
    \end{itemize}

    \column{0.5\textwidth}
    \textbf{Agentic system}
    \begin{itemize}
      \item Iterative control loop
      \item Plan, act, observe, adapt
      \item Tool invocation and validation
      \item Best for multi-step workflows
    \end{itemize}
  \end{columns}
\end{frame}

\begin{frame}{When You Should Not Use an Agent}
  \begin{alertbox}{Common overengineering trap}
    If a deterministic pipeline solves the problem with lower cost and higher reliability, do not use an agent.
  \end{alertbox}

  \vspace{0.4em}
  \begin{itemize}
    \item Static ETL transformations
    \item Low-latency hard real-time tasks
    \item Strictly regulated deterministic outputs
    \item Fixed, non-branching logic
  \end{itemize}
\end{frame}

\begin{frame}{Agent or Workflow? Start from the Control Problem}
  \begin{itemize}
    \item \textbf{Task variability:} Does the workflow require open-ended interpretation or can the branches be enumerated?
    \item \textbf{Action risk:} What is the cost of a wrong tool call, and can the action be reversed or reconciled?
    \item \textbf{Latency and cost:} Is iterative model invocation compatible with the service-level objective and budget?
    \item \textbf{Evidence requirements:} Must each action be justified by auditable rules, retrieved evidence, or human approval?
    \item \textbf{Failure containment:} Can the system stop safely when observations are missing, contradictory, or outside the tested distribution?
  \end{itemize}

  \begin{alertbox}{Design principle}
    Use the least autonomous architecture that satisfies the task. Ambiguity alone is not a reason to grant broad action authority.
  \end{alertbox}
\end{frame}

\begin{frame}{Transition: From Definitions to Design}
  \begin{block}{Next}
    We now move from \textbf{what an agent is} to \textbf{how to structure one} so behavior is predictable and auditable.
  \end{block}
\end{frame}

% ================================================================
\section{Architecture and Orchestration}
% ================================================================

\begin{frame}{Reference Architecture Layers}
  \begin{itemize}
    \item \textbf{Interaction:} request intake, UX, response contracts
    \item \textbf{Policy:} authorization, data governance, action gating
    \item \textbf{Orchestration:} planner, executor, retry, stop logic
    \item \textbf{Model:} reasoning, extraction, ranking, verification
    \item \textbf{Tools:} APIs, DB, workflows, code execution wrappers
    \item \textbf{State:} working memory, session log, long-term knowledge
    \item \textbf{Observability:} traces, metrics, replay artifacts
  \end{itemize}
\end{frame}

\begin{frame}{Control Loop with an Explicit Control Plane}
  \begin{equation*}
    s_{t+1} = T(s_t, a_t, o_t), \qquad a_t = \pi_\theta(s_t)
  \end{equation*}

  \begin{itemize}
    \item $s_t$: recorded application state (goal, plan, constraints, approvals)
    \item $a_t$: proposed action (tool call, clarification, finalize)
    \item $o_t$: observed result, including external-system state and validation outcomes
    \item $T$: application transition logic that should be explicit, bounded, and testable
  \end{itemize}

  \begin{block}{Engineering principle}
    Keep authorization, budgets, stop conditions, and mutation gates in application code. Make transitions deterministic where feasible, but record enough context to diagnose nondeterminism from models, external tools, clocks, and changing data sources.
  \end{block}
\end{frame}

\begin{frame}{Execution Patterns and Trade-offs}
  \begin{table}
    \small
    \begin{tabular}{@{}llll@{}}
      \toprule
      Pattern & Strength & Risk & Use Case \\
      \midrule
      Router & Cheap, simple & Limited depth & Support triage \\
      ReAct & Adaptive & Drift loops & Research/operations \\
      Plan/Execute & Auditable & Planning overhead & Long workflows \\
      DAG workflow & Predictable & Less adaptive & Enterprise flows \\
      Multi-agent & Specialization & Coordination cost & Complex domains \\
      \bottomrule
    \end{tabular}
  \end{table}
\end{frame}

\begin{frame}{Termination Design}
  \textbf{A production agent must know when to stop.}
  \begin{itemize}
    \item Goal achieved with validated outputs
    \item Budget exhausted (tokens, latency, tool-call quota)
    \item Confidence below threshold with no safe next action
    \item Policy gate requires human decision
    \item Repeated failure class without recovery path
  \end{itemize}
\end{frame}

\begin{frame}{Checkpoint 1: Architecture Review}
  \begin{examplebox}{Quick review prompt}
    For your current use case, define:
    \begin{itemize}
      \item Agent objective
      \item Required tools
      \item Termination conditions
      \item Mandatory human approval points
    \end{itemize}
  \end{examplebox}
\end{frame}

% ================================================================
\section{Tools, State, and Memory}
% ================================================================

\begin{frame}[fragile]{Tool Contract Principles}
  \begin{itemize}
    \item Strong input schema
    \item Stable, versioned output envelope
    \item Explicit error taxonomy and retry semantics
    \item Idempotency key or reconciliation strategy for mutating calls
  \end{itemize}

  \begin{lstlisting}[language=Python, caption=Minimal tool envelope]
{
  "ok": true,
  "data": {...},
  "error_code": null,
  "trace_id": "abc-123"
}
  \end{lstlisting}
\end{frame}

\begin{frame}[fragile]{Example: Mutation Tool Contract}
  \begin{lstlisting}[language=Python, caption=Refund action contract]
{
  "name": "create_refund",
  "input_schema": {
    "type": "object",
    "properties": {
      "order_id": {"type": "string"},
      "reason_code": {"type": "string", "enum": ["damaged", "late", "wrong_item"]},
      "amount_cents": {"type": "integer", "minimum": 1}
    },
    "required": ["order_id", "reason_code", "amount_cents"]
  },
  "idempotency_key_required": true
}
  \end{lstlisting}
\end{frame}

\begin{frame}{State Model}
  \textbf{State object should be explicit and minimal.}
  \begin{itemize}
    \item Request metadata and user intent
    \item Active plan and completed steps
    \item Evidence references (not full raw payloads)
    \item Tool execution history with outcomes
    \item Policy flags and pending approvals
  \end{itemize}
\end{frame}

\begin{frame}{Memory Strategy}
  \begin{columns}
    \column{0.5\textwidth}
    \textbf{Working memory}
    \begin{itemize}
      \item Active context
      \item Current constraints
      \item Latest tool outputs
    \end{itemize}

    \textbf{Session memory}
    \begin{itemize}
      \item Prior decisions
      \item User preferences
      \item Interaction history
    \end{itemize}

    \column{0.5\textwidth}
    \textbf{Long-term memory}
    \begin{itemize}
      \item Stable domain facts
      \item Playbooks
      \item Policy documents
    \end{itemize}

    \textbf{Retention and governance}
    \begin{itemize}
      \item Keep credentials and secrets out of model-visible memory
      \item Minimize sensitive payload retention and enforce purpose/TTL/access controls
      \item Preserve only the audit evidence required by policy, regulation, or incident response
    \end{itemize}
  \end{columns}
\end{frame}

\begin{frame}{Retrieval and Grounding}
  \begin{itemize}
    \item Retrieval is evidence support, not authority transfer
    \item Rank by relevance, trust, and freshness
    \item Attach provenance to every grounded claim
    \item Require citation for high-stakes outputs
  \end{itemize}

  \vspace{0.5em}
  \begin{block}{Retrieval evaluation}
    Relevance is only one dimension. Evaluate source authority, freshness, coverage, citation correctness, and whether the retrieved evidence actually supports the downstream claim.
  \end{block}
\end{frame}

\begin{frame}{Checkpoint 2: Tool and Memory Design}
  \begin{examplebox}{Design exercise}
    Pick one high-risk tool in your domain and define:
    \begin{itemize}
      \item Input schema and output envelope
      \item Idempotency strategy
      \item Required approvals
      \item Memory retention policy
    \end{itemize}
  \end{examplebox}
\end{frame}

% ================================================================
\section{Safety, Evaluation, and Operations}
% ================================================================

\begin{frame}{Failure Taxonomy}
  \begin{itemize}
    \item \textbf{Reasoning failures:} invalid decomposition, wrong assumptions
    \item \textbf{Tool failures:} timeouts, schema mismatch, side-effect duplication
    \item \textbf{Grounding failures:} hallucinated or stale claims
    \item \textbf{Policy failures:} unauthorized or unsafe action paths
    \item \textbf{Coordination failures:} state divergence across branches
  \end{itemize}
\end{frame}

\begin{frame}{Practical Safety Controls}
  \begin{itemize}
    \item Pre-action policy checks (role, scope, data class)
    \item Execution sandboxing and network egress controls
    \item Human approval gates for high-impact mutations
    \item Post-action reconciliation checks
    \item Immutable audit logs with trace ids
  \end{itemize}
\end{frame}

\begin{frame}{Evaluation Stack}
  \textbf{Three evaluation levels:}
  \begin{enumerate}
    \item \textbf{Unit-level:} tool contract correctness
    \item \textbf{Workflow-level:} scenario replay and regression sets
    \item \textbf{Production-level:} online metrics and drift monitoring
  \end{enumerate}

  \vspace{0.4em}
  \begin{block}{Stage-appropriate release criteria}
    Pre-production tests should gate known scenarios and safety invariants; shadow/canary stages then add online evidence before broader rollout. Thresholds must be defined for the use case and risk level.
  \end{block}
\end{frame}

\begin{frame}{Core Metrics}
  \begin{columns}
    \column{0.5\textwidth}
    \textbf{Quality}
    \begin{itemize}
      \item Task success rate
      \item Plan validity
      \item Tool-call precision/recall
      \item Unsupported-claim rate under an operational definition
    \end{itemize}

    \column{0.5\textwidth}
    \textbf{Operations}
    \begin{itemize}
      \item p50/p95 latency
      \item Cost per successful task
      \item Retry rate per step
      \item Escalation rate
    \end{itemize}
  \end{columns}
\end{frame}

\begin{frame}{Observability and Debugging}
  \begin{itemize}
    \item Per-step trace with prompt/tool/result metadata
    \item Correlation id across model and tool calls
    \item Structured error events with failure class
    \item Replay or reconstruction mode using recorded model versions, prompts, tool inputs/outputs, and external-state snapshots where available
    \item Diff between expected and observed transitions
  \end{itemize}
\end{frame}

\begin{frame}{Cost Governance}
  \begin{itemize}
    \item Per-request token budgets
    \item Per-agent step limits
    \item Caching for repeated retrieval and deterministic calls
    \item Model routing: expensive models only for critical steps
    \item Cost-per-outcome monitoring, not just token tracking
  \end{itemize}
\end{frame}

\begin{frame}{Transition: From Components to End-to-End Delivery}
  \begin{block}{Next}
    We now combine architecture, tooling, memory, and controls in one production case study.
  \end{block}
\end{frame}

\begin{frame}{Checkpoint 3: Reliability Readiness}
  \begin{examplebox}{Readiness questions}
    \begin{itemize}
      \item Can you reconstruct the important state and evidence for a failed run?
      \item Do you have a regression set for top workflows?
      \item Are high-impact mutations approval-gated and either idempotent or reconcilable?
      \item Do you track cost per successful outcome?
    \end{itemize}
  \end{examplebox}
\end{frame}

% ================================================================
\section{Case Study Walkthrough}
% ================================================================

\begin{frame}{Case Study: Customer Support Resolution Agent}
  \textbf{Goal:} Resolve tickets with minimal human intervention while protecting policy and customer trust.

  \vspace{0.4em}
  \textbf{Tools:}
  \begin{itemize}
    \item CRM lookup
    \item Order and shipment systems
    \item Refund/credit mutation endpoint
    \item Knowledge base retrieval
  \end{itemize}
\end{frame}

\begin{frame}{Case Study Flow}
  \begin{enumerate}
    \item Classify issue and confidence level
    \item Retrieve policy and customer/order evidence
    \item Propose resolution plan with mutation requirements
    \item Run policy checks and approval gate if needed
    \item Execute mutation and verify post-state
    \item Respond with explanation and audit metadata
  \end{enumerate}
\end{frame}

\begin{frame}[fragile]{Case Study: Plan Record}
  \begin{lstlisting}[language=Python, caption=Plan object persisted in state]
{
  "goal": "Resolve ticket #T-90821",
  "steps": [
    {"id": 1, "action": "lookup_order", "status": "done"},
    {"id": 2, "action": "check_policy", "status": "done"},
    {"id": 3, "action": "propose_refund", "status": "done"},
    {"id": 4, "action": "execute_refund", "status": "pending_approval"}
  ],
  "stop_if": ["approval_rejected", "budget_exceeded", "policy_violation"]
}
  \end{lstlisting}
\end{frame}

\begin{frame}{Case Study: Failure and Recovery}
  \textbf{Failure:} refund API timeout after mutation accepted upstream.

  \vspace{0.4em}
  \textbf{Recovery strategy:}
  \begin{itemize}
    \item Retry with same idempotency key
    \item Reconcile state by querying source-of-truth ledger
    \item If inconsistent, escalate with trace id and pause automation
    \item Emit incident signal for reliability dashboard
  \end{itemize}
\end{frame}

\begin{frame}{From Pilot to Production}
  \begin{enumerate}
    \item Start with read-only paths
    \item Introduce one low-risk mutation tool
    \item Add approval gate and post-action verification
    \item Expand coverage by failure-class maturity
    \item Expand only after stage-specific quality, safety, and incident criteria are stable for a period appropriate to traffic volume and risk
  \end{enumerate}
\end{frame}

% ================================================================
\section{Closing}
% ================================================================

\begin{frame}{Anti-Patterns to Avoid}
  \begin{itemize}
    \item Prompt tuning before system boundary definition
    \item Untyped tool calls and implicit side effects
    \item Insufficient trace data to reconstruct failures or compare runs
    \item Missing approval gates on mutating actions
    \item Shipping with no regression benchmark
  \end{itemize}
\end{frame}

\begin{frame}{Production Readiness Checklist}
  \begin{theorembox}{Before you deploy}
    \begin{itemize}
      \item Objective and stop criteria are explicit
      \item Tool interfaces are typed; mutations are idempotent or reconcilable
      \item Transition logic, gates, and recorded state are auditable
      \item Safety controls exist pre-action and post-action
      \item Evaluation covers quality, reliability, latency, and cost
      \item Incident replay is operationally usable
    \end{itemize}
  \end{theorembox}
\end{frame}

\begin{frame}{Synthesis}
  \begin{itemize}
    \item Agent reliability is primarily a systems property: model quality matters, but authorization, tool contracts, state management, stop conditions, and recovery paths determine whether failures are contained.
    \item Autonomy should be scoped to the action risk. Read-only retrieval, reversible changes, and irreversible mutations deserve different control and approval policies.
    \item Evaluation must decompose failure modes. Final-answer quality alone can hide invalid plans, unsupported evidence, unsafe tool calls, or duplicate side effects.
    \item Reproducibility is often partial rather than literal. Record model/version/configuration, prompts, state transitions, retrieved evidence, and tool outcomes so runs can be reconstructed and compared.
    \item More planning, memory, or multiple agents increase capability only when the added coordination cost and failure surface are justified by measured task performance.
  \end{itemize}
\end{frame}

\begin{frame}{Selected References}
  \small
  \begin{itemize}
    \item S. Yao et al. (2023). "ReAct: Synergizing Reasoning and Acting in Language Models." \textit{ICLR}.
    \item T. Schick et al. (2023). "Toolformer: Language Models Can Teach Themselves to Use Tools." \textit{NeurIPS}.
    \item X. Liu et al. (2024). "AgentBench: Evaluating LLMs as Agents." \textit{ICLR}.
    \item Z. Xi et al. (2023). "The Rise and Potential of Large Language Model Based Agents: A Survey." arXiv:2309.07864.
  \end{itemize}
\end{frame}

\begin{frame}{Facilitator Timing Guide}
  \begin{itemize}
    \item \textbf{By slide 10}: finish fundamentals
    \item \textbf{By slide 20}: finish architecture and orchestration
    \item \textbf{By slide 30}: finish tools, memory, retrieval
    \item \textbf{By slide 37}: finish safety/evaluation/operations
    \item \textbf{Slides 38+}: case study, summary, Q\&A
  \end{itemize}
\end{frame}

\begin{frame}{Q\&A}
  \begin{block}{Discussion prompts}
    \begin{itemize}
      \item Which workflow in your team is best for first deployment?
      \item What is your highest-risk mutation action?
      \item Which metric will define success in the first 30 days?
    \end{itemize}
  \end{block}
\end{frame}

\appendix
\section{Appendix: Deep Dives}

\begin{frame}{Deep Dive: Planner Design Choices}
  \begin{itemize}
    \item \textbf{Single-shot planner:} fast, but brittle under changing evidence
    \item \textbf{Rolling planner:} re-plans each step, more robust but costlier
    \item \textbf{Hybrid:} coarse global plan + local rolling re-plan
  \end{itemize}

  \vspace{0.4em}
  \begin{alertbox}{Recommendation}
    Start with hybrid planning for long-running workflows with external dependencies.
  \end{alertbox}
\end{frame}

\begin{frame}{Deep Dive: Tool Error Taxonomy}
  \begin{table}
    \small
    \begin{tabular}{@{}lll@{}}
      \toprule
      Error Class & Retry? & Typical Handling \\
      \midrule
      ValidationError & No & fix input / ask clarification \\
      AuthError & No & escalate / refresh auth context \\
      TimeoutError & Yes & bounded retries + backoff \\
      ConflictError & Maybe & reconcile state, retry if safe \\
      PolicyDenied & No & halt and request approval \\
      \bottomrule
    \end{tabular}
  \end{table}
\end{frame}

\begin{frame}{Deep Dive: Retrieval Evaluation}
  \begin{itemize}
    \item Build query set from real production intents
    \item Measure top-k hit rate and citation correctness
    \item Include freshness tests for time-sensitive policies
    \item Track failure by source tier (internal docs vs external content)
  \end{itemize}
\end{frame}

\begin{frame}{Deep Dive: Regression Harness}
  \begin{itemize}
    \item Store fixed input tasks with expected outcomes
    \item Replay with fixed fixtures or deterministic tool mocks when possible
    \item Compare plan traces, not only final answers
    \item Fail release on safety regressions even if accuracy improves
  \end{itemize}
\end{frame}

\begin{frame}{Deep Dive: Rollout Strategy}
  \begin{enumerate}
    \item Internal alpha with shadow mode
    \item Limited beta with read-only or low-risk actions
    \item Gradual percentage rollout with automatic rollback trigger
    \item Full rollout after stable SLO and incident trend
  \end{enumerate}
\end{frame}

\end{document}
